\documentclass[conference]{IEEEtran}
\usepackage[T1]{fontenc}
\usepackage{booktabs,multirow,amsmath,graphicx,xcolor,url}
\newcommand{\tbd}[1]{\textcolor{red}{[#1]}}
% Double-anonymous review: \anontrue hides the author, the affiliation and the artifact location.
\newif\ifanon
\anonfalse
\newcommand{\confTasks}{30}
\newcommand{\confHeadline}{found the opposite of the exploratory phase: with the same model in every arm, pre-computed
packs cost 41\% less than plain search (95\% CI 34--48\%) and a reduced tool set 35\% less, both at equal or better
quality, and the result held for a smaller model (39\% and 36\%). The full tool set cost 15\% less; task-shaped tools
saved nothing.}
\newcommand{\ablationHeadline}{A 2$\times$2 ablation could not attribute the change to the graph audit or to the
completeness statement: neither moved cost by more than run-to-run noise, while the plain baseline alone varied by
39\% between sessions.}
\newcommand{\processHeadline}{On process-layer joins (ticket to code to tests, commit to epic and specification) packs
cost 60\% less but failed non-inferiority on quality, because two incomplete joins were trusted without checking.}


\begin{document}
\title{Precomputed Context, Not More Tools: A Pre-Registered Study of Process Graphs for Coding Agents}

\ifanon
\author{\IEEEauthorblockN{Anonymous author(s)}\IEEEauthorblockA{Anonymous institution}}
\else
\author{\IEEEauthorblockN{Hung T. Le}\IEEEauthorblockA{Independent researcher, Austin, Texas, USA \quad letronghung@gmail.com}}
\fi

\maketitle

\begin{abstract}
Coding agents pay to rediscover what a repository already records: callers, tests, tickets, specifications. We built
a software-process knowledge graph (20 node types, 30 relations) and served it to agents as tools and as pre-computed
\emph{packs}, with the same model in every arm. In 461 exploratory runs on two industrial repositories, graph-equipped
agents cost 11--24\% more than plain search, through three mechanisms we document. A pre-registered study on
\confTasks{} held-out tasks in four repositories (1,020 runs) found the opposite: packs cost 41\% less (95\% CI
34--48\%) and a small tool set 35\% less, at equal or better quality, for two model sizes; task-shaped tools saved
nothing, and the saving tracked the number of turns. A pre-registered ablation refuted our own explanation that
auditing the graph removed the cost penalty; the plain baseline alone varied 39\% between sessions. On process joins,
packs cost 60\% less but lost quality where the graph had missed a link.
\end{abstract}

\begin{IEEEkeywords}
coding agents, LLM agents, knowledge graphs, software process, retrieval, Model Context Protocol, cost, empirical study, pre-registration
\end{IEEEkeywords}

\section{Introduction}
LLM-based coding agents~\cite{jimenez2024swebench,yang2024sweagent} navigate a repository with
text search and file reads, and every read is paid for again on each later turn. A natural idea is to give the agent
the repository's structure up front: a graph of symbols and calls~\cite{liu2024codexgraph,ouyang2025repograph},
extended with the development process (tickets, commits, specifications, tests, owners, configuration) and served over
a tool protocol~\cite{mcp2024}. For a practitioner the question is not whether such a graph is \emph{useful} but whether
it makes agents \emph{cheaper at equal quality}, on which tasks, and what it takes. Prior work builds such graphs and
the means to navigate them, or measures how well agents find context (Section~\ref{sec:related}); we hold the knowledge
fixed and vary how it is delivered: interactive tools of different sizes, or precomputed context.

We answer with the same model in every arm. Contributions:
\begin{itemize}
\item A software-process knowledge graph for agents, its extraction, and three ways to serve it: single-question
  tools, task tools, and pre-computed \emph{packs} that carry their own evidence and a completeness statement
  (Section~\ref{sec:system}).
\item Three mechanisms by which a graph makes a strong agent \emph{more} expensive, observed in 461 exploratory runs
  (Section~\ref{sec:exploratory}).
\item An edge-level audit that measures each relation against truth computed from source (Section~\ref{sec:audit}).
\item A pre-registered confirmatory study on held-out tasks in four repositories, an ablation separating the audit
  from the completeness statement, and a study of process-layer joins (Sections~\ref{sec:confirmatory}--\ref{sec:process}),
  with an artifact that reproduces every number.
\end{itemize}

\section{The graph and how agents receive it}\label{sec:system}
\paragraph{Schema} Node types: Ticket, Commit, File, Symbol, TestCase, Endpoint, Flow, ErrorSignature, Spec,
Requirement, Evidence, WorkPackage, Milestone, Release, WorkUnit, Component, Subsystem, Warning, Repository and
ConfigKey. Relations include \texttt{CALLS}, \texttt{TESTED\_BY}, \texttt{RAISES}, \texttt{STEP} (ordered call flows
from entry points), \texttt{TOUCHES}, \texttt{REFERENCES}, \texttt{CHILD\_OF}, \texttt{IMPLEMENTS},
\texttt{SIMILAR\_TO}, \texttt{READS} and \texttt{DECLARES}, 30 in all. We call it a knowledge graph rather than an
ontology: it is a typed property graph whose schema changes only through recorded architecture decisions, without
formal semantics.
\paragraph{Extraction} Calls are resolved through each file's imports, aliases, relative imports, typed receivers
(parameter annotations, constructors, \texttt{self} attributes) and base classes; nothing is attributed by a bare name.
Every list answer carries \emph{coverage}: the call sites under the same name the resolver could not attribute, and the
files holding them.
\paragraph{Serving} \emph{Single-question tools} (12): callers, tests for, impact, review pack, ticket context, lineage,
\dots; \emph{task tools} (7): one call per developer task; a \emph{lean} set (7 tools); and \emph{packs}, written
before the session from what a task names, placed in the prompt, with no tools loaded. A pack gives every caller,
handler and test as \texttt{path:line} with the line of code that proves it, a coverage statement, and freshness
against the working tree.

\section{Method}\label{sec:method}
\paragraph{Arms} plain (Read, Grep, Glob, read-only git, a neutral efficiency note); core; workflow; lean; pack. Same
model, same task prompt, a 40-turn cap, no session persistence, a discarded warm-up per arm, runs shuffled with a fixed
seed. Graph arms are told to use the graph first; whether an agent adopts it unprompted is reported separately.
\paragraph{Truth} Computed from source and history, never from the graph: calls resolved by import path, grep, git,
the ticket files' own trace tables.
\paragraph{Measures} API-equivalent cost per run, turns, tool-result bytes, and a task score in $[0,1]$ (F1 of named
sets, exact match of a raiser or a trace). Cost effect $=1-\sum_t \overline{c}_{\mathrm{arm},t}/\sum_t
\overline{c}_{\mathrm{plain},t}$ (positive means cheaper), with 95\% bootstrap intervals resampling runs within task.

\section{Exploratory phase}\label{sec:exploratory}
\begin{table*}[t]\centering\small
\caption{Exploratory benchmarks (Sonnet). ``Saved'' is the cost saved against plain search (positive means cheaper):
E1 reports medians per stage; E2--E5 report per-task bootstrap 95\% intervals. Score is task correctness from 0 to 1,
graph arm against plain search.}
\label{tab:exploratory}
\setlength{\tabcolsep}{5pt}
\begin{tabular}{@{}llrlll@{}}\toprule
Run & Repository and tasks & Runs & Arm & Saved & Score (arm vs.\ plain) \\\midrule
E1 & B, six lifecycle stages & 117 & forced graph, S1 understand & +54\% & 1.00 vs.\ 0.97 \\
   & & & forced graph, S8 review & +53\% & 0.99 vs.\ 0.98 \\
   & & & forced graph, S6 readiness & +61\% & 1.00 vs.\ 1.00 \\
   & & & forced graph, S3 requirements & $-$54\% & 0.60 vs.\ 0.42 \\
   & & & forced graph, S4 solutioning & +7\% & 0.88 vs.\ 0.58 \\\midrule
E2 & B, six developer tasks & 80 & core & $-$11\% [$-$33, +10] & 0.96 vs.\ 1.00 \\
   & & & workflow & $-$15\% [$-$39, +4] & 0.83 vs.\ 1.00 \\\midrule
E3 & A, six multi-hop tasks & 72 & core & $-$9\% [$-$43, +16] & 0.95 vs.\ 0.92 \\
   & & & workflow & $-$6\% [$-$38, +19] & 0.87 vs.\ 0.92 \\\midrule
E4 & A, same tasks, call-resolution fixes & 72 & core & $-$24\% [$-$40, $-$10] & 0.94 vs.\ 0.94 \\
   & & & workflow & $-$11\% [$-$20, $-$2] & 0.91 vs.\ 0.94 \\\midrule
E5 & A, same tasks, audited graph and coverage & 120 & core & +5\% [$-$27, +28] & 0.96 vs.\ 0.93 \\
   & & & lean & +12\% [$-$18, +35] & 0.95 vs.\ 0.93 \\
   & & & pack & +24\% [$-$6, +45] & 0.94 vs.\ 0.93 \\\bottomrule
\end{tabular}
\end{table*}

Table~\ref{tab:exploratory} summarises five exploratory benchmarks. Across the development lifecycle (E1) the graph cut
cost where a stage follows call chains (understanding a function, reviewing a commit, checking readiness) and improved
requirement and solution answers at higher cost. Left to choose, the agent used the graph in 88\% of understanding runs
and in none of the requirement, solution or review runs. On whole developer tasks (E2--E4) no graph arm saved cost; on
Repository A the graph arms cost 11--24\% more, with intervals excluding zero (E4). Three mechanisms recur:
\begin{enumerate}
\item \emph{Per-call overhead.} Tool definitions and replies cost more than the one or two searches they replace for a
  one-hop question.
\item \emph{Re-reading to verify.} A strong model opens the files a tool names before it trusts the answer, so the tool
  adds cost without removing reads.
\item \emph{Complete-looking lists.} When a refactor tool's recall rose from 1 to 11 of 16 caller files, accuracy fell
  from 0.89 to 0.77: a plausible list ended the search, an obviously short one had not.
\end{enumerate}

\section{Auditing the graph against source truth}\label{sec:audit}
\begin{table}[t]\centering\footnotesize
\caption{Repository A's graph before and after the audit, measured against truth computed from source. Precision at 8 of similar tickets is against random ticket pairs at 0.14.}
\label{tab:audit}
\setlength{\tabcolsep}{3pt}
\begin{tabular}{@{}p{0.44\columnwidth}p{0.24\columnwidth}p{0.24\columnwidth}@{}}\toprule
Measure & Before & After \\\midrule
Edges & 352,556 & 167,399 \\
Nodes & 67,049 & 53,740 \\
Size on disk & 109.9 MB & 59.5 MB \\
Call recall / precision & 0.711 / 0.985 & 0.719 / 0.986 \\
Import-based test links: precision (edges) & 1.00 (4,305) & 1.00 (4,305) \\
Ticket-marker test links: precision (edges) & 0.007 (73,132) & 0.01 (11,686) \\
Commit--file edges (most files in one commit) & 134,900 (13,035) & 83,690 (334) \\
Similar tickets per ticket; precision at 8 & 51; 0.47 & 7; 0.47 \\
Frontend-to-route recall & 0 & 0.64 \\
Error and setting-read recall & 1.00 & 1.00 \\
Answers marked complete that are complete & 90.5\% & 100\% \\\bottomrule
\end{tabular}
\end{table}

We measured every relation against source-derived truth (Table~\ref{tab:audit}). Ticket-marker test links (``a test
names ticket $T$, so it tests every file $T$ touched'') had 0.7\% precision on 73k edges and were read by no tool;
38 bulk commits (under 1\% of commits) produced 40\% of commit--file edges; ticket similarity kept 51 neighbours per
ticket where readers used 8; and the frontend-to-backend relation was empty because API calls built from template
literals and wrappers were skipped. Restricting, capping or repairing these cut edges by 53\% and disk size by 46\%,
kept call precision at 0.986, and made every answer marked complete actually complete (from 90.5\%). After the audit
and the coverage statement, the graph arms' cost effect moved from $-24\%$ to $+5\%$ (E4 to E5). We first read this
as the audit's effect; the ablation (Section~\ref{sec:ablation}) does not support that reading.

\section{Confirmatory study}\label{sec:confirmatory}
The exploratory phase tuned tools on its own tasks, so we froze the system and a protocol (hypotheses, outcomes,
analysis, sample size) before generating any new task. Tasks were generated with a new seed in four repositories
(Table~\ref{tab:repos}), excluding every function used before. Five deviations are listed in the artifact. Two came
before any run: call truth resolves by import path instead of requiring unique names (D1), and one-hop error traces are
admitted (D2). 763 runs that the account's usage limit stopped before they started were re-run in full cells (D3).
Two concern the analysis: how the no-worse hypothesis H3 is tested (D4), and an ablation baseline from a later
session (D5).
\begin{table}[t]\centering\footnotesize
\caption{Repositories of the confirmatory study, each pinned at one commit.}
\label{tab:repos}
\begin{tabular}{@{}llrrr@{}}\toprule
Id & Repository & Py files & Routes & Tasks \\\midrule
A & industrial cybersecurity platform & 957 & 242 & 7 \\
B & industrial software-factory tooling & 375 & 72 & 8 \\
C & Netflix Dispatch (Apache-2.0) & 717 & 290 & 8 \\
D & Mealie (AGPL-3.0) & 720 & 291 & 7 \\\bottomrule
\end{tabular}
\end{table}

\begin{table}[t]\centering\small
\caption{Confirmatory study: cost saved against plain (positive is cheaper) and quality difference, 95\% bootstrap intervals; 30 held-out tasks, 750 Sonnet and 270 Haiku runs. Plain: median 6 turns (Sonnet), mean score 0.83.}\label{tab:confirmatory}
\setlength{\tabcolsep}{3pt}
\resizebox{\columnwidth}{!}{\begin{tabular}{llrrrl}\toprule
Model & Arm & Cost saved & Quality diff. & Turns & Hyp. \\ \midrule
Sonnet & core & +15\% [+7, +22] & +0.023 [+0.007, +0.039] & 7 & H3 \\
Sonnet & workflow & $-$4\% [$-$13, +4] & +0.045 [+0.029, +0.061] & 9 &  \\
Sonnet & lean & +35\% [+28, +42] & +0.053 [+0.036, +0.072] & 6 & H2 \\
Sonnet & pack & +41\% [+34, +48] & +0.034 [+0.019, +0.048] & 3 & H1 \\
\midrule
Haiku & lean & +36\% [+31, +41] & +0.056 [+0.024, +0.088] & 7 & H5 \\
Haiku & pack & +39\% [+31, +48] & +0.038 [+0.007, +0.071] & 4 & H5 \\
\bottomrule\end{tabular}}
\end{table}

% Numbers from results-sonnet.json and results-haiku.json (analyze.py); table from fill.py.
All 1,020 runs produced an answer that could be scored (one Sonnet pack run was invalid and scored 0). Within each
repository every arm ran in the same session in shuffled order (Deviation D3 re-ran whole repositories, never single
arms).

\paragraph{Hypotheses} \emph{H1 supported}: packs cost 41\% less than plain [34, 48], Holm-adjusted $p<10^{-4}$, with
quality $+0.034$ [0.019, 0.048]. \emph{H2 supported}: the lean tool set cost 35\% less [28, 42], $p<10^{-4}$, quality
$+0.053$. \emph{H3 supported as stated}: the full tool set cost no more than plain; it cost 15\% less [7, 22] with
quality $+0.023$. Read as a superiority claim inside the Holm family it would not reach significance (adjusted
$p=0.13$); the pre-registration did not say which reading applies, so we report both (Deviation D4). \emph{H4 not
supported}: only two error-to-handler tasks qualified, and a one-sided Wilcoxon test on two tasks cannot go below
$p=0.25$ although both arms were cheaper on both tasks (pack 65\%, lean 33\%). \emph{H5 supported}: with Haiku, packs
cost 39\% less [31, 48] and lean 36\% less [31, 41], both $p<10^{-3}$ and quality non-inferior.

\paragraph{Where the saving comes from} The saving followed the turn count: packs answered in a median of 3 turns
against 6 for plain, and the task-shaped tools, which took 9, saved nothing ($-4\%$ [$-13$, 4]). By task kind, packs
saved 46\% on change-impact tasks (15) and 26\% on refactor tasks (13); the full tool set saved 28\% on change-impact
and cost 16\% \emph{more} on refactors, where callers must be complete. By repository, pack savings ranged from 15\%
(C and D, the public repositories; D's interval includes zero) to 55--68\% (A, B). Quality was non-inferior pooled,
but not in every repository: on Repository A, packs scored 0.029 lower [$-0.054$, $-0.009$].

\paragraph{Why the sign flipped} The exploratory phase found graph arms 11--24\% more expensive, the confirmatory study
15--41\% cheaper. Three things changed between them, and this design cannot separate them: the graph was audited and
answers gained a coverage statement (Section~\ref{sec:audit}); the tasks are new and held out; and the tools were
frozen after tuning on the exploratory tasks. Section~\ref{sec:ablation} tests the first.


\section{Ablation: content or completeness statement?}\label{sec:ablation}
% Numbers from ablation.json (analyze_ablation.py).
The ablation crosses graph content (G0: the graph before the audit; G1: after) with the coverage statement (C0: off;
C1: on), using the full tool set on the six exploratory tasks of Repository A, four runs per cell (96 runs, plus 24
plain).

\paragraph{Result} Neither factor moved cost measurably. The main effect of the audit on the cost saving was $-1.2$
points [$-7.5$, 5.3], of the coverage statement $+4.4$ [$-2.2$, 10.9], and their interaction $-2.3$ [$-15.2$, 10.7].
Quality was 0.93--0.96 in every cell. All four cells, the unaudited graph without coverage among them, cost 22--28\%
less than the plain baseline.

\paragraph{A baseline from another session} That last comparison should not be trusted, and the reason is a finding in
itself. The plain baseline was re-run in a later session after the usage limit stopped it (Deviation D3), while the
four cells ran together. On the same six tasks and the same model, the plain arm's total mean cost was \$0.57, \$0.41
and \$0.55 in three sessions: a 39\% swing with no change to the arm. Contrasts between cells that ran together are
sound; a contrast against an arm from another session is not (Deviation D5).

\paragraph{Consequence} The exploratory change from $-24\%$ to $+5\%$ (E4 to E5) lies within the between-session spread
of the plain arm, and the ablation finds no effect of the audit or of the coverage statement on cost at this size. We
therefore withdraw the exploratory reading that the audit removed the cost penalty. The audit's value is in what the
graph says (precision, honest coverage), not demonstrably in what the agent spends.


\section{Process-layer joins}\label{sec:process}
% Numbers from results-process.json (analyze_process.py).
Sixteen tasks in Repositories A and B (eight per repository): \emph{ticket-to-code} (which top-level functions did a
ticket's commits change that still exist, and which test files call them; truth from git history and import
resolution) and \emph{commit-to-trace} (a commit's ticket, that ticket's epic and specification; truth from the
ticket's own trace table). Arms plain, core and pack, five runs each (240 runs, all valid).

\paragraph{H6 not supported} Packs cost 60\% less than plain [48, 70] ($p<10^{-4}$), answering in one turn against
seven, but scored 0.080 lower [$-0.094$, $-0.058$], outside the non-inferiority margin. Plain search was nearly perfect
on these tasks (mean score 0.99). The full tool set was no cheaper ($-7\%$ [$-23$, 7]) and scored 0.12 lower.

\paragraph{Where it failed} On commit-to-trace, packs matched plain exactly (score 1.0) at 37\% less cost. On
ticket-to-code, the loss came from two of eight tasks. In one, the graph did not link the ticket to its commits; in
the other, the pack's test list missed most of the test files. Each time the agent accepted the pack in one turn and
did not search. This is the third exploratory mechanism, complete-looking lists, recurring in a pre-registered
setting. The ticket pack's coverage statement reports call sites it could not attribute, but not commits it failed
to link to the ticket, so a missed join looked like a complete answer. Coverage has to extend to every join a pack
makes.


\section{Discussion}
\paragraph{Precomputed context, not more tools} On held-out tasks the saving tracked the number of turns: packs (median 3
turns) and a small tool set (6) saved 35--41\%, task-shaped tools (9) saved nothing. What a graph buys an agent is the
removal of turns, and every tool definition and every reply it has to verify gives some of that back.
\paragraph{Joins are where packs are trusted most and checked least} On process joins, packs were cheapest and the
agent did not check them; two incomplete joins cost the study its quality margin. A pack must state what it could not
join, not only what it could not resolve.
\paragraph{Mechanisms outlast percentages} Prices, caching and models' habits change; the three mechanisms of
Section~\ref{sec:exploratory} are properties of how agents use tools, and two of them reappeared in the confirmatory
studies (task tools that cost turns; complete-looking joins).
\paragraph{Compare arms only within a session} The plain arm alone varied by 39\% between sessions on identical tasks.
An exploratory before-and-after comparison of that size is noise; the confirmatory design ran every arm of a
repository in one shuffled session, and we recommend the same to anyone benchmarking agent tooling.

\section{Threats to validity}
\emph{Construct}: scores compare against an import-resolved truth that does not see calls through untyped objects;
tasks were selected where that truth is complete. \emph{Internal}: the author built tools, tasks and truth; mitigated by
pre-registration, held-out tasks, two public repositories and a full artifact. The ablation and the process study are
secondary and pre-registered in an addendum. \emph{External}: Python web systems, two model sizes from one vendor,
\confTasks{} confirmatory tasks; error-to-handler tasks were rare in the held-out repositories, so that hypothesis is
underpowered; process tasks come from the two internal repositories only, because the public ones keep their tickets
outside the repository. Pack savings were smaller on the two public repositories (15\%) than on the internal ones
(55--68\%), so the size of the effect depends on the codebase, and the percentages are results on these tasks, not a
general saving for coding agents: on one-hop questions the saving was small. \emph{Conclusion}: per-task bootstrap and Holm
correction; score differences below about 0.05 per task kind are within run-to-run noise; cost varies by up to 39\%
between sessions, so only within-session contrasts are interpreted (Deviations D3--D5).

\section{Related work}\label{sec:related}
\emph{Repository graphs for agents.} RepoGraph~\cite{ouyang2025repograph} adds a graph of definitions and references
to existing agent frameworks; CodexGraph~\cite{liu2024codexgraph} puts a code graph database behind a query-writing
agent; LocAgent~\cite{chen2025locagent} gives agents a few traversal tools over a code graph for localization. These
works build the graph and the means to navigate it; we fix a graph, extend it to the development process, and vary only
how its answers reach the agent.
\emph{Repository context.} ContextBench~\cite{li2026contextbench}, SWE-Explore~\cite{zhang2026sweexplore} and Agent
Retrieval Bench~\cite{qin2026arb} measure how well agents or retrievers find the context a task needs. SWE-Explore is
closest in design, handing an explorer's output to a fixed agent; we hand the agent precomputed answers from a graph and
compare them with interactive tools over the same graph, by cost at equal quality.
\emph{Interfaces and cost.} The agent-computer interface shapes results~\cite{yang2024sweagent}, a fixed pipeline can
beat complex agents at lower cost~\cite{xia2025agentless}, agents should be reported on cost and accuracy
together~\cite{kapoor2025agents}, and many tools degrade tool selection~\cite{gan2025ragmcp}; SWE-bench~\cite{jimenez2024swebench}
measures resolution, not cost at equal quality.

\section{Conclusion}
On held-out tasks, with the same model in every arm, a software-process knowledge graph made coding agents 35--41\%
cheaper at equal quality when it was delivered as a pre-computed pack or a small tool set, for a large and a small
model. Task-shaped tools saved nothing, and on process joins the cheapest delivery lost quality where the graph had
missed a link. Our own exploratory explanation, that auditing the graph removed a cost penalty, did not survive the
ablation. The practical rule is the one the data support: deliver the graph in as few turns as possible, and make
every answer say what it could not see.

\paragraph*{AI use and artifact} Per IEEE policy: an AI coding assistant helped implement the extractors, tools and
harness, run the benchmarks, and draft the text; the author designed the study, checked every result against the raw
data, and is responsible for the content. \ifanon An anonymised artifact (code, tasks with truth, raw runs, analysis,
pre-registration) is available to reviewers.\else Code, protocol, analysis scripts and raw runs:
\url{https://github.com/namiq-ai/devkb} (industrial repositories as per-run metrics only).\fi

\bibliographystyle{IEEEtran}
\bibliography{refs}
\end{document}
